% !TEX root = ../../Main.tex
\section{Motivation}

Automated currency trading is usually built from fixed rules. Technical indicators are computed from recent prices. A signal is then taken from them, either by comparing them or by a crossover. A separate layer of the system decides the position size \cite{chan_quantitative_2021, brandimarte_introduction_2018}. These systems are cheap to run and easy to check, and this is why they are common. Their parameters are set once and do not change, so the system does not adapt when the market changes. This is the first of the two difficulties that this work deals with.

Currencies were chosen over equities for practical reasons \cite{donnelly_art_2019}. The market trades continuously on weekdays, so a position is never carried through a daily close during which it cannot be traded. There are only seven major pairs. This makes it possible to test a design on all of them, instead of on a sample chosen by the author. A trader can buy or sell with no borrowing cost and no short-selling restrictions. A policy that sells is therefore not penalised by market mechanics that have nothing to do with its signal. Every major pair also has the US dollar on one side. The seven instruments are therefore related, but one cannot replace another. A comparison of the results across them says something about the design and not about one price series.

Reinforcement learning describes the same task without splitting it into a prediction and an action. The agent observes the state of the market and chooses a position. It then receives the profit or loss of that position as its reward \cite{sutton_reinforcement_2018}. The agent never has to forecast a price. It only has to act. Its action is measured in money, which is also how a trading system is judged in practice.

DDPG fits this task well. It is an actor-critic method for continuous action spaces \cite{lillicrap_continuous_2015}. The output of the actor can therefore be used directly as a signed position. The sign gives the direction, and the magnitude gives the exposure. Much of the published work on trading uses value-based methods such as the Deep Q-Network \cite{mnih_human-level_2015, fischer_reinforcement_2018, hambly_recent_2023}. These methods need the action to be one of a small discrete set, normally buy, hold and sell. Position sizing is then decided outside the learned policy, by rules that were never trained on the reward.

The second difficulty is cost. It does not depend on the learning algorithm. Results in this field are often reported on bar data, with the spread fixed or ignored and without the broker commission \cite{fischer_reinforcement_2018, pricope_deep_2021}. For a strategy that trades a few times per year, this changes little. For a strategy that changes its position often, it changes a lot. The cost is paid on every change, while the profit per trade stays small. A system that makes money without costs can lose money with them. The test only shows this if it charges what the account would pay \cite{lopez_de_prado_advances_2018}. The size of this difference is easy to underestimate. In this work, charging the broker commission on an otherwise identical run removed twenty percentage points of total return from a configuration that changed its position at every decision. Once a threshold was added so that small changes were skipped, it removed six percentage points. Both figures are reported in \Cref{Chapter:ExperimentalResults}.